% ---------------------------------------------------------------------------
% Shared style for the Hybroid classifier dot plots (thesis Figures 4.7, 4.8).
%
% Colours are taken verbatim from the TUM Corporate Design print palette:
%   #0065BD  TUM Blau, Pantone 300 C, used at 100 per cent (no tints of the
%            brand colour are permitted), reserved here for the fused modality
%   #64A0C8  Akzentfarbe, Pantone 542
%   #98C6EA  Akzentfarbe, Pantone 283
%   #CCCCCC  20 per cent black
%   #808080  50 per cent black
%   #333333  80 per cent black
% The three series form a lightness ramp rather than a hue contrast, so the
% figure survives greyscale printing and red-green colour deficiency. Orange
% and green are absent, as the CD requires when marks distinguish categories.
%
% Type is TeX Gyre Heros, the Helvetica clone already loaded by the thesis
% style, so figure labels match the document's sans-serif face.
% ---------------------------------------------------------------------------
\usepackage[T1]{fontenc}
\usepackage{tgheros}
\renewcommand{\familydefault}{\sfdefault}
\usepackage{tikz}
\usetikzlibrary{calc}

\definecolor{tumblue}{HTML}{0065BD}
\definecolor{tumacc542}{HTML}{64A0C8}
\definecolor{tumacc283}{HTML}{98C6EA}
\definecolor{tumgreyxx}{HTML}{CCCCCC}
\definecolor{tumgreyl}{HTML}{808080}
\definecolor{tumgreyd}{HTML}{333333}

% Geometry, in millimetres. The total width matches \textwidth of the thesis
% (402.565 pt = 141.6 mm), so the figure is included at scale 1:1 and its type
% sizes are the sizes that reach the page.
\def\LX{20}        % width of the shared learner-label column
\def\PW{37.2}      % plot width of one metric panel
\def\GAP{5}        % gap between panels
% The score axis is set per figure, because the two tasks occupy different
% ranges and a shared axis would leave half of the detection panels empty.
% Each caption states the range it uses. \TICKS lists the gridded, labelled
% positions; both figures step by 0.04 so the two look alike.
\def\VMIN{0.78}
\def\VMAX{0.98}
\def\TICKS{0.80,0.85,0.90,0.95}
\def\ROWA{-7}
\def\ROWB{-15}
\def\ROWC{-23}
\def\AXY{-27.5}    % baseline of the score axis
\def\GTOP{-3.6}    % top of the grid lines
\def\GBOT{-26.4}   % bottom of the grid lines

% Map a score onto an x coordinate inside panel #1 (0, 1 or 2).
\newcommand{\px}[2]{\LX+#1*(\PW+\GAP)+(#2-\VMIN)/(\VMAX-\VMIN)*\PW}

\newcommand{\ticklabelfont}{\fontsize{6.2}{7.4}\selectfont}
\newcommand{\rowlabelfont}{\fontsize{7}{8.4}\selectfont}
\newcommand{\paneltitlefont}{\fontsize{7.6}{9}\selectfont\bfseries}

% One metric panel: grid, axis, tick labels, title.
\newcommand{\panel}[2]{%
  \foreach \v in \TICKS{%
    \draw[tumgreyxx,line width=0.3pt] ({\px{#1}{\v}},\GTOP) -- ({\px{#1}{\v}},\GBOT);
    \node[below,text=tumgreyl,inner sep=1.1pt] at ({\px{#1}{\v}},\AXY)
      {\ticklabelfont\v};
  }%
  \draw[tumgreyd,line width=0.4pt] ({\px{#1}{\VMIN}},\AXY) -- ({\px{#1}{\VMAX}},\AXY);
  \node[above right,text=black,inner sep=0pt] at ({\px{#1}{\VMIN}},1.2)
    {\paneltitlefont #2};
}

% Markers. Program code and network flow are peer modalities, so they carry
% the same radius and are separated by lightness alone; the fused marker is
% larger and takes the brand colour, because it is the result the chapter
% argues about. Drawing runs largest first, so where a single modality and the
% fused representation report the same score the two show as concentric discs
% rather than one hiding the other. Neither figure contains a program
% code / network flow tie, which is the one case this ordering cannot resolve.
\newcommand{\mcode}[2]{\draw[fill=tumacc283,draw=tumacc542,line width=0.35pt]
  ({#1},#2) circle (0.58mm);}
\newcommand{\mflow}[2]{\draw[fill=tumacc542,draw=tumacc542,line width=0.3pt]
  ({#1},#2) circle (0.58mm);}
\newcommand{\mfused}[2]{\draw[fill=tumblue,draw=tumblue,line width=0.3pt]
  ({#1},#2) circle (0.88mm);}

% One row of a panel: connector spanning the observed range, then the three
% markers from largest to smallest. #1 panel, #2 row y, #3 program code,
% #4 network flow, #5 fused.
\newcommand{\dotrow}[5]{%
  \pgfmathsetmacro{\vlo}{min(#3,min(#4,#5))}%
  \pgfmathsetmacro{\vhi}{max(#3,max(#4,#5))}%
  \draw[tumgreyxx,line width=0.7pt] ({\px{#1}{\vlo}},#2) -- ({\px{#1}{\vhi}},#2);
  \mfused{\px{#1}{#5}}{#2}\mflow{\px{#1}{#4}}{#2}\mcode{\px{#1}{#3}}{#2}%
}

% Learner labels, written once for all three panels.
\newcommand{\rowlabels}{%
  \node[left,inner sep=0pt] at (\LX-2.4,\ROWA) {\rowlabelfont Decision tree};
  \node[left,inner sep=0pt] at (\LX-2.4,\ROWB) {\rowlabelfont Random forest};
  \node[left,inner sep=0pt] at (\LX-2.4,\ROWC) {\rowlabelfont Gradient boosting};
}

% Legend, one row under the panels.
\newcommand{\legendrow}[1]{%
  \def\ly{#1}%
  \mcode{\LX}{\ly}
  \node[right,inner sep=0pt] at (\LX+1.4,\ly) {\rowlabelfont Program code};
  \mflow{\LX+26}{\ly}
  \node[right,inner sep=0pt] at (\LX+27.4,\ly) {\rowlabelfont Network flow};
  \mfused{\LX+52}{\ly}
  \node[right,inner sep=0pt] at (\LX+53.4,\ly) {\rowlabelfont Fused representation};
}
